%% =============================================================
%% Section 2: Model
%% =============================================================

\section{Model}
\label{sec:model}

\subsection{Environment}
\label{sec:environment}

Time is discrete with two periods, $t \in \{1,2\}$. The economy is a small open
economy populated by a representative household-firm unit. The central bank sets the
nominal interest rate $i_t$ each period to minimise a quadratic loss function. We use
lower-case variables to denote log-deviations from steady state; all steady-state
values are normalised to zero.

\paragraph{IS curve.}
Aggregate demand in each period follows the standard open-economy Euler equation:
\begin{equation}
  y_t = \mathbb{E}_t[y_{t+1}] - \sigma\bigl(i_t - \mathbb{E}_t[\pi_{t+1}]\bigr)
        + \alpha\, q_t + \varepsilon_{d,t},
  \label{eq:IS}
\end{equation}
where $y_t$ is the output gap, $\pi_t$ is CPI inflation, $i_t$ is the nominal policy
rate, $q_t$ is the (log) real exchange rate (an increase is a depreciation), and
$\varepsilon_{d,t}$ is a demand disturbance. The parameter $\sigma > 0$ is the
intertemporal elasticity of substitution and $\alpha > 0$ captures the expenditure-
switching effect of the real exchange rate on net exports.

\paragraph{New Keynesian Phillips curve.}
Inflation dynamics are governed by
\begin{equation}
  \pi_t = \beta\, \mathbb{E}_t[\pi_{t+1}] + \kappa\, y_t + \varepsilon_{c,t},
  \label{eq:NKPC}
\end{equation}
where $\beta \in (0,1)$ is the household discount factor, $\kappa > 0$ is the slope
of the Phillips curve, and $\varepsilon_{c,t}$ is a cost-push shock.

\paragraph{Uncovered interest parity.}
The real exchange rate satisfies the forward-looking UIP condition:
\begin{equation}
  q_t = \mathbb{E}_t[q_{t+1}] - \bigl(i_t - i^*\bigr),
  \label{eq:UIP}
\end{equation}
where $i^* \geq 0$ is the (exogenous, constant) foreign interest rate. Because the
model has two periods, the exchange rate today embeds the entire expected path of
future interest rate differentials. Iterating~\eqref{eq:UIP} forward and imposing the
terminal condition $\mathbb{E}_2[q_3] = 0$ gives:
\begin{equation}
  q_1 = -\bigl(i_1 - i^*\bigr) - \mathbb{E}_1\bigl[i_2 - i^*\bigr].
  \label{eq:q1_path}
\end{equation}
Equation~\eqref{eq:q1_path} is central to the paper's mechanism: the current exchange
rate is pinned by the entire expected path of the policy rate. Any uncertainty about
$i_2$ therefore transmits immediately into $q_1$.

\paragraph{Terminal conditions.}
Period 2 is the final period. We impose:
\begin{equation}
  \mathbb{E}_2[y_3] = \mathbb{E}_2[\pi_3] = \mathbb{E}_2[q_3] = 0.
  \label{eq:terminal}
\end{equation}
Substituting~\eqref{eq:terminal} into~\eqref{eq:IS}--\eqref{eq:NKPC} yields the
period-2 block:
\begin{align}
  y_2 &= -\sigma\, i_2 + \alpha\, q_2 + \varepsilon_{d,2}, \label{eq:IS2}\\
  \pi_2 &= \kappa\, y_2 + \varepsilon_{c,2}. \label{eq:NKPC2}
\end{align}
From~\eqref{eq:UIP} and the terminal condition:
\begin{equation}
  q_2 = -\bigl(i_2 - i^*\bigr).
  \label{eq:q2}
\end{equation}
Substituting~\eqref{eq:q2} into~\eqref{eq:IS2} and defining the \emph{effective
interest rate sensitivity}
\begin{equation}
  \tilde{\sigma} \equiv \sigma + \alpha > \sigma,
  \label{eq:sigma_tilde}
\end{equation}
the period-2 output gap (after normalising $\alpha i^* = 0$) becomes:
\begin{equation}
  y_2 = -\tilde{\sigma}\, i_2 + \varepsilon_{d,2}.
  \label{eq:IS2_reduced}
\end{equation}
The parameter $\tilde{\sigma}$ exceeds the closed-economy elasticity $\sigma$ because a
rate increase both suppresses domestic spending \emph{and} appreciates the exchange
rate, reducing net exports.

%% -----------------------------------------------------------
\subsection{Central Bank Preferences and Type}
\label{sec:CB_prefs}

The central bank minimises the standard quadratic loss function:
\begin{equation}
  \mathcal{L} = \sum_{t=1}^{2} \beta^{t-1}
    \Bigl[ y_t^2 + \lambda\,\pi_t^2 \Bigr],
  \label{eq:loss}
\end{equation}
where $\lambda > 0$ is the relative weight on inflation stabilisation. A higher
$\lambda$ represents a more \emph{hawkish} central bank; a lower $\lambda$ a more
\emph{dovish} one.

\paragraph{Private information.}
The inflation weight $\lambda$ is the central bank's \emph{private information}. The
private sector holds a prior over two possible types:
\begin{equation}
  \lambda \in \{\lambda_L,\, \lambda_H\}, \quad \lambda_L < \lambda_H,
  \label{eq:types}
\end{equation}
with prior probabilities
\begin{equation}
  \Pr(\lambda = \lambda_H) = p \in (0,1), \qquad
  \Pr(\lambda = \lambda_L) = 1-p.
  \label{eq:prior}
\end{equation}
Private agents know the distribution~\eqref{eq:prior} and the model equations but
cannot directly observe $\lambda$.

\paragraph{Remark.}
The assumption that $\lambda$ is private information is motivated by empirical evidence
that central bank preferences are not perfectly observable to the public: markets spend
considerable resources inferring the degree of inflation aversion from policy
decisions, communications, and appointments. This uncertainty about the reaction
function---rather than uncertainty about economic fundamentals---is the source of
welfare losses in the opacity regime.

%% -----------------------------------------------------------
\subsection{Timing and Communication Regimes}
\label{sec:timing}

The sequence of events is as follows:

\begin{enumerate}[label=\arabic*.]
  \item \textbf{Period 1 (announcement stage).} The cost-push shock $\varepsilon_{c}$
        is publicly observed.\footnote{This captures the realistic scenario in which
        the central bank and private sector alike observe current inflationary
        pressure---for example, an energy price spike---before the bank announces its
        policy intentions.} The central bank sets $i_1$. Under the two regimes:
        \begin{itemize}
          \item \textit{Forward guidance (FG):} The central bank announces
                $\hat{\imath}_2$, a forecast of its intended period-2 rate. The
                announcement is truthful and believed in full by private agents.
          \item \textit{Opacity (OP):} The central bank makes no announcement.
                Private agents form $\mathbb{E}_1^{\mathrm{OP}}[i_2]$ using
                Bayes' rule over the type distribution.
        \end{itemize}
  \item \textbf{Period 2 (implementation stage).} The demand shock $\varepsilon_{d}$
        is realised and observed by all. The central bank sets $i_2$ by minimising its
        period-2 loss, taking the realised shocks and its own type as given. There is
        \emph{no commitment}: the bank is free to deviate from any period-1
        announcement.
\end{enumerate}

Under \emph{forward guidance}, the announcement $\hat{\imath}_2$ is the central bank's
rational forecast of its own period-2 decision, formed in period 1 before $\varepsilon_d$
is observed:
\begin{equation}
  \hat{\imath}_2(\lambda) \equiv \mathbb{E}_1\bigl[i_2^*(\lambda,
  \varepsilon_d, \varepsilon_c)\bigr]
  = \phi_c(\lambda)\,\varepsilon_c,
  \label{eq:FG_announcement}
\end{equation}
where the second equality follows from $\mathbb{E}_1[\varepsilon_d] = 0$ and the
period-2 optimal policy derived below (Proposition~\ref{prop:period2}). The
announcement is a forecast, not a promise: the bank will re-optimise in period 2 upon
observing $\varepsilon_d$.

%% -----------------------------------------------------------
\subsection{Period-2 Equilibrium}
\label{sec:period2}

\begin{proposition}[Period-2 Optimal Policy]
\label{prop:period2}
In period 2, the central bank of type $\lambda$ minimises $y_2^2 + \lambda\pi_2^2$
subject to \eqref{eq:IS2_reduced} and \eqref{eq:NKPC2}. The optimal policy rate is:
\begin{equation}
  i_2^*(\lambda) = \phi_d\,\varepsilon_{d,2} + \phi_c(\lambda)\,\varepsilon_{c},
  \label{eq:i2star}
\end{equation}
where
\begin{align}
  \phi_d &\equiv \frac{1}{\tilde{\sigma}},
  \label{eq:phi_d}\\[4pt]
  \phi_c(\lambda) &\equiv
    \frac{\lambda\kappa}{\tilde{\sigma}\bigl(1 + \lambda\kappa^2\bigr)}.
  \label{eq:phi_c}
\end{align}
\end{proposition}

\begin{proof}
Substituting~\eqref{eq:IS2_reduced} into~\eqref{eq:NKPC2}:
\begin{equation*}
  \pi_2 = -\kappa\tilde{\sigma}\,i_2 + \kappa\varepsilon_{d,2} + \varepsilon_c.
\end{equation*}
The first-order condition $\partial(y_2^2 + \lambda\pi_2^2)/\partial i_2 = 0$ gives the
optimal trade-off condition:
\begin{equation}
  y_2 + \lambda\kappa\,\pi_2 = 0.
  \label{eq:FOC_period2}
\end{equation}
Substituting $y_2 = -\tilde{\sigma}i_2 + \varepsilon_d$ and
$\pi_2 = -\kappa\tilde{\sigma}i_2 + \kappa\varepsilon_d + \varepsilon_c$ into
\eqref{eq:FOC_period2} and solving for $i_2$ yields~\eqref{eq:i2star}. \qed
\end{proof}

\paragraph{Interpretation of the response coefficients.}
Several properties of Proposition~\ref{prop:period2} deserve emphasis.

\begin{enumerate}[label=(\roman*)]
  \item \textit{Demand shocks are accommodated identically by all types.}
        The coefficient $\phi_d = 1/\tilde{\sigma}$ is independent of $\lambda$.
        Because a demand shock raises output and inflation proportionally, it creates
        no tension between stabilisation objectives; both types respond by raising the
        rate to fully close the output gap.

  \item \textit{Cost-push shocks reveal the central bank's type.}
        The coefficient $\phi_c(\lambda)$ is strictly increasing in $\lambda$:
        \begin{equation}
          \frac{\partial\phi_c}{\partial\lambda}
            = \frac{\kappa}{\tilde{\sigma}(1+\lambda\kappa^2)^2} > 0.
          \label{eq:dphi_dlambda}
        \end{equation}
        A hawkish central bank ($\lambda_H$) raises the rate more aggressively in
        response to a cost-push shock than a dovish one ($\lambda_L$), accepting a
        larger output contraction to contain inflation. It is therefore the
        \emph{response to cost-push shocks} that identifies the type.

  \item \textit{The exchange rate amplifies the policy transmission.}
        The effective sensitivity $\tilde{\sigma} = \sigma + \alpha > \sigma$ means
        that, for a given change in $i_2$, output and inflation move more than in a
        closed economy. This amplification is internalised by the central bank: a
        smaller rate move is needed to achieve the same stabilisation, so $i_2^*$ is
        lower in absolute value than its closed-economy counterpart.
\end{enumerate}

\paragraph{Period-2 equilibrium outcomes.}
Substituting~\eqref{eq:i2star} back into \eqref{eq:IS2_reduced} and \eqref{eq:NKPC2}:
\begin{align}
  y_2^*(\lambda) &= \frac{-\lambda\kappa}{1+\lambda\kappa^2}
    \bigl(\kappa\varepsilon_{d,2} + \varepsilon_c\bigr)
    + \frac{\varepsilon_{d,2}}{1 + \lambda\kappa^2} \cdot \frac{1}{1},
  \label{eq:y2star}
\end{align}
more compactly:
\begin{align}
  y_2^*(\lambda)   &= \frac{\varepsilon_{d,2} - \lambda\kappa\varepsilon_c}{1+\lambda\kappa^2},
  \label{eq:y2star_clean}\\[4pt]
  \pi_2^*(\lambda) &= \frac{\kappa\varepsilon_{d,2} + \varepsilon_c}{1+\lambda\kappa^2}.
  \label{eq:pi2star_clean}
\end{align}

\noindent Note from \eqref{eq:pi2star_clean} that the denominator $(1+\lambda\kappa^2)$
is increasing in $\lambda$: a hawkish central bank achieves smaller inflation
deviations in response to cost-push shocks, but at the cost of larger output gap
deviations (equation~\eqref{eq:y2star_clean}).

%% -----------------------------------------------------------
\subsection{Period-1 Equilibrium}
\label{sec:period1}

We now characterise the period-1 equilibrium under each communication regime. The
central bank sets $i_1$ to minimise total expected loss~\eqref{eq:loss}, taking as
given the private sector's expectations (which differ across regimes) and the
period-2 equilibrium~\eqref{eq:i2star}.

\paragraph{Expectations under each regime.}
The period-2 rate is $i_2^*(\lambda) = \phi_d\varepsilon_d + \phi_c(\lambda)\varepsilon_c$.
Since $\mathbb{E}_1[\varepsilon_d] = 0$, the critical object is how agents form
$\mathbb{E}_1[i_2]$:

\begin{itemize}
  \item \textbf{Forward guidance:} The announcement $\hat{\imath}_2(\lambda) =
        \phi_c(\lambda)\varepsilon_c$ is observed and believed. Thus:
        \begin{equation}
          \mathbb{E}_1^{\mathrm{FG}}[i_2] = \phi_c(\lambda)\,\varepsilon_c.
          \label{eq:E_i2_FG}
        \end{equation}

  \item \textbf{Opacity:} Agents do not observe $\lambda$. Using the prior
        \eqref{eq:prior} and the known equilibrium mapping~\eqref{eq:i2star}:
        \begin{equation}
          \mathbb{E}_1^{\mathrm{OP}}[i_2]
            = \bigl[p\,\phi_c(\lambda_H) + (1-p)\,\phi_c(\lambda_L)\bigr]\varepsilon_c
            \equiv \bar{\phi}_c\,\varepsilon_c,
          \label{eq:E_i2_OP}
        \end{equation}
        where $\bar{\phi}_c \equiv p\,\phi_c(\lambda_H)+(1-p)\,\phi_c(\lambda_L)$
        is the prior-weighted average response coefficient.
\end{itemize}

\paragraph{Exchange rate in period 1.}
From~\eqref{eq:q1_path} and normalising $i^* = 0$:
\begin{equation}
  q_1^r = -i_1 - \mathbb{E}_1^r[i_2], \qquad r \in \{\mathrm{FG}, \mathrm{OP}\}.
  \label{eq:q1_regime}
\end{equation}
The exchange rate gap between regimes is therefore:
\begin{equation}
  q_1^{\mathrm{FG}} - q_1^{\mathrm{OP}}
    = -\bigl(\phi_c(\lambda) - \bar{\phi}_c\bigr)\varepsilon_c.
  \label{eq:q1_gap}
\end{equation}
For a hawkish central bank ($\lambda = \lambda_H$), we have $\phi_c(\lambda_H) >
\bar{\phi}_c$, so $q_1^{\mathrm{FG}} < q_1^{\mathrm{OP}}$: the announcement of a
hawkish type \emph{appreciates} the exchange rate relative to the opacity benchmark.
The opposite holds for a dovish type. Under opacity, the exchange rate is anchored to
the prior-weighted average response---a systematic misalignment whenever $\lambda \neq
\bar{\lambda}$ (where $\bar{\lambda} = p\lambda_H + (1-p)\lambda_L$).

\paragraph{Period-1 IS curve.}
We substitute $\mathbb{E}_1[y_2]$, $\mathbb{E}_1[\pi_2]$, and $q_1$ into
\eqref{eq:IS}. Using the period-2 equilibrium:
\begin{align}
  \mathbb{E}_1[y_2]  &= 0, \label{eq:E1y2}\\
  \mathbb{E}_1[\pi_2] &= \frac{\kappa\varepsilon_c}{1+\lambda\kappa^2}
    \quad (\text{FG}), \qquad
    \mathbb{E}_1[\pi_2] = \kappa\,\mathbb{E}\!\left[
      \frac{\varepsilon_c}{1+\lambda\kappa^2}\right]
    \quad (\text{OP}),
  \label{eq:E1pi2}
\end{align}
where in \eqref{eq:E1y2} we used $\mathbb{E}[y_2^*] = 0$ under both regimes (since
shocks are mean-zero from the period-1 perspective, using $\mathbb{E}[\varepsilon_d] = 0$
and the structure of $y_2^*$).

Substituting into the period-1 IS curve:
\begin{equation}
  y_1^r = -\sigma\bigl(i_1 - \mathbb{E}_1^r[\pi_2]\bigr)
           -\alpha\bigl(i_1 + \mathbb{E}_1^r[i_2]\bigr)
         = -\tilde{\sigma}\,i_1
           + \sigma\,\mathbb{E}_1^r[\pi_2]
           - \alpha\,\mathbb{E}_1^r[i_2].
  \label{eq:IS1_reduced}
\end{equation}

\paragraph{Period-1 NKPC.}
\begin{equation}
  \pi_1^r = \beta\,\mathbb{E}_1^r[\pi_2] + \kappa\,y_1^r.
  \label{eq:NKPC1_r}
\end{equation}

Equations~\eqref{eq:IS1_reduced}--\eqref{eq:NKPC1_r} define the period-1 allocation
under each regime as a function of $i_1$ and the regime-specific expectation
$\mathbb{E}_1^r[i_2]$. The central bank chooses $i_1$ to minimise expected total loss,
taking~\eqref{eq:i2star} as the period-2 continuation.
